\section{Database Management System (DBMS)}

\begin{definitionbox}{DBMS (ডিবিএমএস)}
ডেটাবেজ তৈরি এবং তার উপাত্ত সংরক্ষণ, অনুসন্ধান, সংশোধন ও ব্যবহার নিয়ন্ত্রণের সফটওয়্যারকে Database Management System (DBMS) বলে।
\end{definitionbox}

\subsubsection{অনলাইন অর্ডার ব্যবস্থাপনায় DBMS}
একটি অনলাইন দোকানে ক্রেতা অর্ডার দিচ্ছেন, কর্মী ঠিকানা সংশোধন করছেন এবং হিসাব বিভাগ টাকা পরিশোধের অবস্থা দেখছে। প্রত্যেকে আলাদা ফাইলের কপি বদলালে কোন তথ্যটি সর্বশেষ, তা বোঝা কঠিন হয়। কে কোন তথ্য বদলাতে পারবেন, সেটিও নিয়ন্ত্রণ করতে হয়।

অ্যাপ্লিকেশন DBMS-এর মাধ্যমে ডেটা পড়লে ও পরিবর্তন করলে একই জায়গা থেকে প্রয়োজনীয় নিয়ম ও অনুমতি প্রয়োগ করা যায়। একই ব্যবস্থা কলেজের ভর্তি, ব্যাংকের হিসাব ও হাসপাতালের নিবন্ধনেও কাজে লাগে। ওয়েবপৃষ্ঠার বোতাম কাজের অনুরোধ পাঠায়; পেছনে ডেটা পরিচালনা করে DBMS।

Microsoft Access, SQLite, MySQL ও Oracle Database পরিচিত DBMS-এর উদাহরণ। 



\subsubsection{Metadata ও Data Dictionary}
\begin{definitionbox}{}
\textbf{Metadata} হলো ডেটার কাঠামো ও নিয়ম সম্পর্কে বর্ণনামূলক তথ্য।
\end{definitionbox}
StudentID-এর টাইপ, কোন কলাম Primary Key এবং GPA-এর গ্রহণযোগ্য সীমা মেটাডেটার উদাহরণ। এগুলো শিক্ষার্থীর প্রকৃত রেকর্ড নয়; রেকর্ড কীভাবে গঠিত হবে, তার বর্ণনা। DBMS-এর \textbf{Data Dictionary বা System Catalog}-এ এসব বর্ণনা থাকে। DBMS এই বর্ণনা ব্যবহার করে নির্ধারিত নিয়ম অনুযায়ী ডেটা গ্রহণ ও ব্যবহার করে।

\subsubsection{DBMS-এর বৈশিষ্ট্য}
\begin{enumerate}
\item \textbf{Physical Data Independence:} যৌক্তিক কাঠামো অপরিবর্তিত রেখে ভেতরের সংরক্ষণপদ্ধতি পরিবর্তনের সুযোগ দেয়।
\item \textbf{একাধিক View:} একই মূল ডেটা থেকে ব্যবহারকারীর কাজ অনুযায়ী ভিন্ন উপস্থাপন তৈরি করা যায়।
\item \textbf{অনুমতি নিয়ন্ত্রণ:} কে কোন ডেটা দেখতে বা পরিবর্তন করতে পারবেন, তা নিয়ন্ত্রণ করে।
\end{enumerate}
\textbf{উদাহরণ:} অনুসন্ধানের জন্য Index যোগ করলেও Name কলামের অর্থ বদলায় না। অফিসের কর্মী ফি এবং শিক্ষক ফল দেখতে পারেন; প্রত্যেকের জন্য সব তথ্যের পৃথক কপি লাগে না।
\textbf{মনে রাখবে:} Physical Data Independence মানে যেকোনো কাঠামোগত পরিবর্তনের পর সব অ্যাপ্লিকেশন স্বয়ংক্রিয়ভাবে ঠিক কাজ করবে—এমন নিশ্চয়তা নয়।

\subsection{Database System ও কাজের প্রবাহ}
ডেটাবেজ, DBMS, হার্ডওয়্যার, ব্যবহারকারী, অ্যাপ্লিকেশন সফটওয়্যার ও পরিচালনার নিয়ম নিয়ে একটি \textbf{Database System} গঠিত হয়। ব্যবহারকারী ফর্ম বা নির্দেশ দিয়ে কাজের অনুরোধ করেন। DBMS তাঁর অনুমতি ও ডেটার নিয়ম পরীক্ষা করে প্রয়োজনীয় কাজ সম্পন্ন করে এবং ফল দেয়। ব্যবহারকারীকে ডিস্কে প্রতিটি রেকর্ডের অবস্থান সরাসরি নিয়ন্ত্রণ করতে হয় না।

ব্যবহারকারীর অনুরোধ ও ফলের সম্পর্ক নিচের কার্যপ্রবাহে দেখানো হয়েছে। Front End ও Back End-এর পৃথক ভূমিকা পরের উপাংশে আলোচনা করা হয়েছে।


\begin{flushleft}
\begin{tikzpicture}[node distance=4mm,every node/.style={booknode,font=\small,text width=24mm,minimum height=11mm}]
\node (u) {User / App};
\node[right=of u] (m) {DBMS};
\node[right=of m] (d) {Database};
\node[right=of d] (r) {Result};
\draw[bookarrow] (u)--(m);\draw[bookarrow] (m)--(d);\draw[bookarrow] (d)--(r);
\end{tikzpicture}
\end{flushleft}
এটি কাজের ধারণাগত প্রবাহ; ফলও DBMS ও অ্যাপ্লিকেশনের মাধ্যমে ব্যবহারকারীর কাছে আসে।
\subsection{DBMS-এর প্রধান কাজ}
\textbf{কাঠামো তৈরি:} DBMS দিয়ে ডেটাবেজ ও টেবিল তৈরি, ফিল্ডের নাম ও টাইপ নির্ধারণ এবং কী ও বিভিন্ন বিধিনিষেধ বসানো যায়। যেমন GPA ফিল্ড সংখ্যাসূচক রাখা এবং 0 থেকে 5-এর মধ্যে মান গ্রহণের নিয়ম নির্ধারণ করা যায়। প্রয়োজনে পরে নতুন ফিল্ড যোগ বা কাঠামো সংশোধন করা যায়।

\textbf{ডেটা ব্যবহার:} নতুন রেকর্ড সংযোজন, প্রয়োজনীয় রেকর্ড অনুসন্ধান, মান পরিবর্তন, রেকর্ড মুছে ফেলা এবং শর্ত অনুযায়ী ফল বের করা DBMS-এর মৌলিক কাজ। পরীক্ষায় GPA 5 পাওয়া শিক্ষার্থীদের তালিকা, বিভাগের মোট শিক্ষার্থী বা বকেয়া ফি-এর হিসাব এভাবে পাওয়া যায়। ফর্ম ও রিপোর্ট ব্যবহার করে এন্ট্রি এবং উপস্থাপন সহজ করা হয়।

\textbf{রক্ষণাবেক্ষণ ও নিয়ন্ত্রণ:} ব্যবহারকারীর পরিচয় ও অনুমতি পরীক্ষা, একাধিক ব্যবহারকারীর কাজের সমন্বয়, ডেটার সামঞ্জস্য বজায় রাখা, ব্যাকআপ ও পুনরুদ্ধারে সহায়তা করা DBMS-এর গুরুত্বপূর্ণ দায়িত্ব। একই আসন দুই গ্রাহকের নামে একসঙ্গে বরাদ্দ না হওয়া প্রয়োজন; এ ধরনের লেনদেন সঠিকভাবে সম্পন্ন করার জন্য নিয়ন্ত্রণের ব্যবস্থা থাকে।

\begin{flushleft}\small
\begin{tabularx}{\linewidth}{|p{32mm}|X|}\hline
\chSixHeader কাজ & কলেজের উদাহরণ\\\hline
কাঠামো সংজ্ঞায়ন & StudentID, Name, BirthDate দিয়ে টেবিল তৈরি\\\hline
সংযোজন ও সংশোধন & নতুন ভর্তি এবং যোগাযোগের নম্বর হালনাগাদ\\\hline
অনুসন্ধান ও বিশ্লেষণ & নির্দিষ্ট বিভাগের ফল ও গড় নম্বর বের করা\\\hline
নিয়ন্ত্রণ & শিক্ষককে নম্বর সংশোধনের অনুমতি, শিক্ষার্থীকে নিজ ফল দেখার অনুমতি\\\hline
সংরক্ষণ ও পুনরুদ্ধার & নিয়মিত ব্যাকআপ রাখা এবং ত্রুটির পর তথ্য ফিরিয়ে আনা\\\hline
\end{tabularx}
\end{flushleft}


\textbf{একই ভর্তি ব্যবস্থায় কাজগুলো চেনো।} ভর্তি ব্যবস্থার টেবিল তৈরি এবং ডেটা রাখার নিয়ম নির্ধারণ করা
কাঠামো নির্ধারণের কাজ। এ কাজের নির্দেশসমষ্টিকে \textbf{Data Definition Language বা DDL} বলে।
নতুন শিক্ষার্থীর রেকর্ড যোগ করা বা তার ফোন নম্বর সংশোধন করা
ডেটা পরিবর্তনের কাজ। এ কাজের নির্দেশসমষ্টিকে \textbf{Data Manipulation Language বা DML} বলে। নির্দিষ্ট বিভাগের শিক্ষার্থীদের তালিকা বের করা তথ্য আহরণের কাজ (Retrieval)।

আবার শিক্ষককে শুধু নিজের বিষয়ের নম্বর পরিবর্তনের অধিকার দেওয়া ব্যবহার নিয়ন্ত্রণের কাজ (Access Control)।
\textbf{Structured Query Language (SQL)}-এর নির্দেশের বিস্তারিত শ্রেণিবিভাগ কুয়েরি অংশে আলোচনা করা হয়েছে; এখানে প্রতিটি কাজের উদ্দেশ্য বোঝাই প্রধান।

কম্পিউটারের আলাদা স্প্রেডশিট বা ফাইলেও তথ্য রাখা যায়; \textbf{ফাইলভিত্তিক ব্যবস্থা মানেই কাগজের খাতা নয়}। পৃথক ফাইলের নকশা ও যাচাই বিভিন্ন প্রোগ্রামে ছড়িয়ে থাকলে একই ফোন নম্বর এক ফাইলে বদলে অন্যটিতে পুরোনো থাকতে পারে। DBMS-ভিত্তিক সমন্বিত নকশায় একটি নির্ভরযোগ্য উৎস ব্যবহার সহজ হয়; তবে সফটওয়্যার বসালেই সব ফাইল স্বয়ংক্রিয়ভাবে মিলে যায় না।

\subsection{Validation, Integrity ও Redundancy}
নির্ধারিত নিয়মের সঙ্গে ইনপুট মিলিয়ে দেখাকে \textbf{Data Validation} বলে। নম্বর 0--100 এবং জন্মতারিখ গ্রহণযোগ্য সীমায় থাকার নিয়ম দিয়ে অবৈধ মান প্রতিরোধ করা যায়। \textbf{Data Integrity} বলতে ডেটার নির্ভুলতা, বৈধতা ও সামঞ্জস্য বজায় থাকার গুণ বোঝায়। অনন্য পরিচিতি এবং অন্য টেবিলে অস্তিত্বশীল রেকর্ডকে উল্লেখ করার নিয়মও ইন্টেগ্রিটি রক্ষা করে।

একই তথ্য অপ্রয়োজনে বহু স্থানে রাখাকে \textbf{Data Redundancy} বলে। একটি কপি সংশোধন করে অন্যগুলো না করলে তথ্যের মধ্যে অসামঞ্জস্য বা \textbf{Data Inconsistency} দেখা দেয়। শিক্ষার্থীর ঠিকানা প্রতিটি পরীক্ষার রেকর্ডে না রেখে পরিচয় টেবিলে একবার রাখা যায়। ফলাফলের টেবিলে তাঁর ID রাখলে প্রয়োজনের সময় পরিচয় ও ফল যুক্ত করা যায়। সংযোগের জন্য ID পুনরায় থাকা প্রয়োজনীয়; সব পুনরাবৃত্ত মানই নকশার ত্রুটি নয়।

\subsection{সুবিধা ও সীমাবদ্ধতা}
\noindent\begin{minipage}[t]{0.485\linewidth}\vspace{0pt}
\begin{keypointbox}{সুবিধা}
\begin{enumerate}
\item \textbf{অনুসন্ধান:} সুষ্ঠু DBMS ব্যবহারে তথ্য দ্রুত পাওয়া সহজ হয়।
\item \textbf{সমন্বিত ব্যবহার:} একাধিক কাজে একই ডেটা ব্যবহার সহজ হয়।
\item \textbf{অনুমতি নিয়ন্ত্রণ:} ডেটা ব্যবহারের অনুমতি নিয়ন্ত্রণ সহজ হয়।
\item \textbf{প্রতিবেদন:} নিয়মিত প্রতিবেদন তৈরি সহজ হয়।
\item পরিকল্পিত নকশা অপ্রয়োজনীয় পুনরাবৃত্তি ও বিরোধপূর্ণ তথ্য কমায়।
\item তথ্যের যৌক্তিক কাঠামো ও ভৌত সংরক্ষণ আলাদা রাখার ফলে কিছু পরিবর্তন অ্যাপ্লিকেশনে কম প্রভাব ফেলতে পারে; এটি ডেটা স্বাধীনতার (Data Independence) ধারণার সঙ্গে সম্পর্কিত।
\end{enumerate}
\end{keypointbox}\end{minipage}\hfill
\begin{minipage}[t]{0.485\linewidth}\vspace{0pt}\begin{mistakebox}{সীমাবদ্ধতা}
\begin{enumerate}
\item \textbf{ব্যয়:} সফটওয়্যার, হার্ডওয়্যার ও প্রশিক্ষণে খরচ হয়.
\item \textbf{জটিল পরিচালনা:} বড় সিস্টেমের নকশা ও রক্ষণাবেক্ষণ জটিল; ভুল অনুমতি বা দুর্বল নিরাপত্তায় বহু তথ্য একসঙ্গে ক্ষতিগ্রস্ত হতে পারে।
\item \textbf{কাজে বিঘ্ন:} কেন্দ্রীয় সিস্টেম অচল হলে অনেক ব্যবহারকারীর কাজ আটকে যায়।
\item \textbf{দক্ষ পরিচালনার প্রয়োজন:} DBMS থাকলেই সব তথ্য নির্ভুল ও নিরাপদ থাকবে না; দক্ষ পরিচালনা, যাচাই ও পুনরুদ্ধারের পরিকল্পনা প্রয়োজন।
\end{enumerate}
\end{mistakebox}\end{minipage}

\subsection{Database Administrator (DBA)}
\begin{definitionbox}{Database Administrator — DBA}
ডেটাবেজ পরিকল্পনা, পরিচালনা ও নিয়ন্ত্রণের দায়িত্বপ্রাপ্ত ব্যক্তিকে \textbf{ডেটাবেজ অ্যাডমিনিস্ট্রেটর বা DBA} বলে।
\end{definitionbox}
DBMS হলো সফটওয়্যার, আর তা পরিচালনার দায়িত্বপ্রাপ্ত DBA একজন ব্যক্তি।


\subsubsection{DBA-এর দায়িত্ব}
\begin{enumerate}
\item \textbf{ব্যবহারকারী ও অনুমতি:} ব্যবহারকারী ও তাঁদের কাজের অনুমতি নির্ধারণ করেন।
\item \textbf{নকশা ও সংরক্ষণ:} নকশা ও সংরক্ষণব্যবস্থা তদারক করেন।
\item \textbf{কার্যক্ষমতা:} সিস্টেমের কার্যক্ষমতা পর্যবেক্ষণ করেন।
\item \textbf{ব্যাকআপ ও পুনরুদ্ধার:} ব্যাকআপ-পুনরুদ্ধারের ব্যবস্থা পরিচালনা করেন।
\end{enumerate}
DBA কেবল ডেটা এন্ট্রি করেন এমন কর্মচারী নন।

\textbf{দায়িত্বের প্রয়োগ:}
\textbf{কার্যক্ষমতা ও ব্যাকআপ:} পরীক্ষার ফল প্রকাশের সময় সিস্টেমের ওপর চাপ বাড়লে DBA সিস্টেমের কার্যক্ষমতা দেখেন, প্রয়োজনীয় ইনডেক্স বা কুয়েরির ব্যবহার পর্যালোচনা করেন এবং ব্যাকআপ যাচাই করেন।

\textbf{অনুমতি নিয়ন্ত্রণ:} নতুন শিক্ষক যোগ হলে তাঁর কাজ অনুযায়ী অনুমতি দেন; শিক্ষক দায়িত্ব ছাড়লে সেই অধিকার প্রত্যাহারের ব্যবস্থা করেন।

\textbf{পরিবর্তনের প্রভাব:} নকশা পরিবর্তনের আগে বিদ্যমান তথ্য ও অ্যাপ্লিকেশনে প্রভাব বিবেচনা করেন।

ডেটা এন্ট্রি অপারেটরের দায়িত্ব সঠিক নথি থেকে তথ্য লেখা; শিক্ষক তাঁর বিষয়ের নম্বর অনুমোদন করেন; DBA কারিগরি ও প্রশাসনিক নিয়ন্ত্রণ পরিচালনা করেন। প্রত্যেকের দায়িত্ব আলাদা থাকলে ভুলের কারণ খুঁজে পাওয়া এবং অনিয়ন্ত্রিত পরিবর্তন প্রতিরোধ করা সহজ হয়।

\subsection{Front End ও Back End-এর ভূমিকা}
একটি ভর্তি ব্যবস্থায় কর্মী পর্দায় আবেদনকারীর নাম লিখে Save দেন। Front End হলো ব্যবহারকারীর দৃশ্যমান interface—পর্দার ঘর, বোতাম ও বার্তা এর অংশ। Back End-এ অনুরোধ প্রক্রিয়াকরণ ও data-access logic থাকে; DBMS ও ডেটাবেজ সেই ব্যবস্থার অংশ হতে পারে। Back End শুধু DBMS-এর আরেক নাম নয়। অনুরোধটি গ্রহণ করে অনুমতি যাচাই, বৈধ মান পরীক্ষা এবং টেবিলে রেকর্ড সংরক্ষণের কাজ পেছনের সেবা ও DBMS করে। ফ্রন্ট এন্ডে ত্রুটি ধরা সুবিধাজনক হলেও শুধু পর্দার যাচাইয়ের ওপর নির্ভর করলে অন্য পথে পাঠানো ভুল ডেটা ঢুকতে পারে। তাই উপযুক্ত বিধি ডেটাবেজ পর্যায়েও থাকতে হয়।
\begin{flushleft}\small
\begin{tabularx}{\linewidth}{|p{30mm}|X|X|}\hline
\chSixHeader বিবেচনা & ফ্রন্ট এন্ড & ব্যাক এন্ড\\\hline
ব্যবহারকারীর সামনে & ফর্ম, অনুসন্ধান পর্দা, ফলের দৃশ্য & সংরক্ষণ ও পরিচালনার অন্তর্গত অংশ\\\hline
উদাহরণ & নাম লিখে Save বোতাম চাপা & INSERT গ্রহণ, বিধি পরীক্ষা, রেকর্ড সংরক্ষণ\\\hline
নিয়ন্ত্রণ & ব্যবহার সহজ করা, প্রাথমিক যাচাই & অধিকার ও ডেটার সামঞ্জস্য প্রয়োগ\\\hline
\end{tabularx}
\end{flushleft}
একটি DBMS একই প্রতিষ্ঠানের বিভিন্ন ফ্রন্ট এন্ডের অনুরোধ গ্রহণ করতে পারে। কলেজের ওয়েবসাইট, অফিসের সফটওয়্যার ও রিপোর্টের অ্যাপ্লিকেশন একই তথ্য ব্যবহার করতে পারে। সব অ্যাপ্লিকেশনকে একই নিয়ম মানাতে কেন্দ্রীয় নকশা ও নিয়ন্ত্রণের গুরুত্ব থাকে।

প্রতিষ্ঠানভিত্তিক ব্যবহার overview, Corporate Database ও সরকারি প্রয়োগে আছে। অতিরিক্ত পরিবহন, টেলিযোগাযোগ, উৎপাদন ও মানবসম্পদের উদাহরণ সংরক্ষিত প্রয়োগ অংশে রাখা হয়েছে।

\subsection{গুরুত্বপূর্ণ প্রশ্ন}
\textbf{অনুশীলন:} (১) ভর্তি ফর্ম ও সংরক্ষিত টেবিলের মধ্যে কোনটি ফ্রন্ট এন্ডের অংশ? (২) 85-এর বদলে 58 লিখলে 0--100 সীমার নিয়ম কেন ভুলটি ধরবে না? (৩) একজন শিক্ষার্থীকে সব শিক্ষার্থীর ফল মুছে ফেলার অধিকার না দেওয়ার দায়িত্ব কার?\par



